\chapter{Local anaesthetic}
\index{Local anaesthetic}

\section*{Definition and Overview}
A local anaesthetic is a pharmacological agent used to prevent or relieve pain in a specific, localised area of the body by reversibly blocking nerve conduction. Unlike general anaesthetics, local anaesthetics do not depress consciousness or affect systemic physiological systems when administered in therapeutic doses. They function by binding to and blocking voltage-gated sodium channels on nerve membranes, which prevents the influx of sodium ions required for nerve depolarisation and the transmission of pain signals to the brain. Local anaesthetics are widely used in clinical practice for minor surgical procedures, dental work, wound repair, pain management during childbirth (such as epidural blocks), and diagnostic interventions. They can be administered via topical application, local infiltration, or regional nerve blocks.

\section*{Epidemiology}
Local anaesthetics are among the most frequently administered classes of drugs worldwide, utilized daily across general practice, dentistry, emergency medicine, and surgery. In many countries, millions of doses are administered annually. The incidence of serious adverse drug reactions is very low (less than 1\%), and true IgE-mediated allergic reactions are extremely rare, accounting for less than 1\% of all reported reactions; most adverse events are instead related to psychogenic responses or accidental intravascular injection. Local Anaesthetic Systemic Toxicity (LAST) is a rare but life-threatening complication, occurring in approximately 0.03\% to 0.1\% of peripheral nerve blocks and epidural procedures. The risk of toxicity is higher in paediatric and geriatric cohorts, pregnant patients, and those with underlying cardiac or hepatic impairment.

\section*{Aetiology and Pathophysiology}
The primary mechanism of local anaesthetics is the reversible inhibition of action potential propagation along peripheral and central nerve fibres. The chemical structure of these agents dictates their metabolism, duration of action, and systemic toxicity profile.

\begin{itemize}
    \item \textbf{Mechanism of action:} Local anaesthetics diffuse across the neuronal sheath in their un-ionised form and bind to the intracellular pore of voltage-gated sodium channels. This blockade prevents sodium influx, stopping depolarisation and halting the generation and propagation of pain impulses.
    \item \textbf{Chemical classification:} They are divided into two main classes: amides (e.g., lignocaine, bupivacaine, ropivacaine, prilocaine) which are metabolised in the liver, and esters (e.g., cocaine, procaine, tetracaine) which are rapidly hydrolysed by plasma pseudocholinesterases. Amide-type agents are more commonly used today.
    \item \textbf{Fibre susceptibility:} Nerve block occurs in a predictable sequence based on fibre size and myelination: small, unmyelinated C-fibres (mediating pain and temperature) are blocked first, followed by small myelinated A-delta fibres, and finally larger myelinated A-alpha and A-beta fibres (mediating touch, pressure, and motor function).
    \item \textbf{Vasoactive properties:} Most local anaesthetics (with the exceptions of cocaine and ropivacaine) cause vasodilation, which increases systemic absorption. For this reason, vasoconstrictors such as adrenaline are frequently added to reduce local bleeding, prolong the duration of the anaesthetic effect, and minimise systemic toxicity.
\end{itemize}

\section*{Clinical Features}
The clinical presentation of local anaesthetic use includes the expected sensory and motor blockade, as well as features of toxicity if the drug enters the systemic circulation in excessive amounts.

\begin{itemize}
    \item \textbf{Sensory and motor blockade:} Loss of pain sensation, followed by loss of temperature, touch, and eventually motor control in the target area.
    \item \textbf{Neurological toxicity (LAST):} Early signs of systemic toxicity include perioral numbness, a metallic taste in the mouth, tinnitus, slurred speech, visual disturbances, lightheadedness, and muscle twitching. This can rapidly progress to generalized tonic-clonic seizures, obtundation, and coma.
    \item \textbf{Cardiovascular toxicity (LAST):} Marked bradycardia, conduction blocks, hypotension, ventricular arrhythmias (such as ventricular fibrillation, particularly associated with bupivacaine), and myocardial depression leading to cardiovascular collapse.
    \item \textbf{Allergic reactions:} Localised pruritus, urticaria, erythema, and bronchospasm. Ester-type agents are more likely to cause allergic reactions due to their metabolite, para-aminobenzoic acid (PABA).
    \item \textbf{Methaemoglobinaemia:} Cyanosis, fatigue, dyspnoea, and chocolate-coloured blood, specifically caused by the accumulation of ortho-toluidine (a metabolite of prilocaine or benzocaine) which oxidises haemoglobin and impairs oxygen delivery.
\end{itemize}

\section*{Differential Diagnosis}
When a patient experiences an adverse reaction following the administration of a local anaesthetic, clinicians must differentiate toxicity from other common procedural events.

\begin{itemize}
    \item \textbf{Vasovagal syncope:} A common psychogenic response to needles or pain. It presents with bradycardia, hypotension, pallor, sweating, and lightheadedness, which resolve rapidly when the patient is placed in a supine position, unlike systemic toxicity.
    \item \textbf{Adrenaline reaction:} Caused by the systemic absorption of co-administered adrenaline. It is characterised by transient tachycardia, hypertension, palpitations, anxiety, and tremors, without the neurological symptoms of local anaesthetic toxicity.
    \item \textbf{Anaphylaxis:} A severe, systemic allergic reaction characterised by bronchospasm, laryngeal oedema, hypotension, and hives, which requires immediate adrenaline administration.
    \item \textbf{Hyperventilation syndrome:} Anxiety-induced tachypnoea causing respiratory alkalosis, which presents with lightheadedness, digital spasm, and perioral paraesthesia, but lacks tinnitus or a metallic taste.
\end{itemize}

\section*{When to Seek Medical Care}
Local anaesthetics are typically administered in clinical settings where patients are monitored. However, if a patient has been discharged and develops delayed symptoms or signs of a complication, they must seek prompt medical care.

\textbf{Contact your local emergency services for:}
\begin{itemize}
    \item Difficulty breathing, wheezing, or swelling of the face, lips, tongue, or throat.
    \item Muscle twitching, tremors, or seizures.
    \item Severe dizziness, confusion, slurred speech, or loss of consciousness.
    \item Palpitations, chest tightness, or a sensation of a racing or irregular heartbeat.
    \item A bluish discolouration of the lips, skin, or fingernails.
\end{itemize}

\section*{Diagnosis and Investigations}
The diagnosis of local anaesthetic toxicity is clinical and requires immediate action. Investigations are used to support management and identify predisposing factors.

\begin{itemize}
    \item \textbf{Clinical monitoring:} Continuous observation of vital signs, level of consciousness, and cardiac rhythm during regional or high-dose blocks.
    \item \textbf{Electrocardiography (ECG):} Used to identify widening of the QRS complex, PR prolongation, bradyarrhythmias, or ventricular arrhythmias associated with cardiotoxicity.
    \item \textbf{Arterial blood gas (ABG) and co-oximetry:} Performed if methaemoglobinaemia is suspected to measure oxygen saturation and direct methaemoglobin levels.
    \item \textbf{Allergy testing:} Referral to an immunologist for skin prick or intradermal testing with preservative-free local anaesthetics if a true IgE-mediated allergy is suspected.
\end{itemize}

\section*{Management}
The management of local anaesthetic administration involves ensuring proper dosing and technique, while treating any complications immediately.

\begin{itemize}
    \item \textbf{Immediate cessation:} Stop the administration of the local anaesthetic at the first sign of systemic toxicity or allergic reaction.
    \item \textbf{Airway management:} Administer 100\% oxygen and ensure adequate ventilation to prevent hypoxia and hypercapnia, which exacerbate local anaesthetic toxicity.
    \item \textbf{Lipid emulsion therapy:} The cornerstone of treatment for severe systemic toxicity (LAST). Intravenous administration of a 20\% lipid emulsion (e.g., Intralipid) acts as a "lipid sink" to extract lipid-soluble anaesthetics from the myocardium and brain tissue.
    \item \textbf{Seizure control:} Intravenous benzodiazepines (e.g., midazolam or diazepam) or propofol are administered to terminate seizure activity.
    \item \textbf{Cardiovascular support:} Advanced Life Support (ALS) guidelines are followed for cardiac arrest, with prolonged resuscitation efforts as local anaesthetic clearing takes time. Avoid bupivacaine for arrhythmia management.
    \item \textbf{Treatment of methaemoglobinaemia:} Intravenous administration of methylene blue is the specific antidote for prilocaine- or benzocaine-induced methaemoglobinaemia.
\end{itemize}

\section*{Complications}
While generally safe, local anaesthetics can be associated with localized and systemic complications.

\begin{itemize}
    \item \textbf{Local Anaesthetic Systemic Toxicity (LAST):} Neurological and cardiovascular collapse from accidental intravascular injection or rapid systemic absorption.
    \item \textbf{Nerve injury:} Direct mechanical trauma from the needle, chemical toxicity of the agent, or ischaemia of the nerve sheath, leading to persistent paraesthesia or motor weakness.
    \item \textbf{Surgical site complications:} Local haematoma, bruising, infection, or tissue necrosis (especially if adrenaline is used in terminal vascular beds like fingers, toes, or the nose).
    \item \textbf{Methaemoglobinaemia:} Impaired oxygen carrying capacity of the blood, which can be life-threatening if untreated.
    \item \textbf{Self-inflicted injury:} Biting the tongue, lip, or cheek, or sustaining burns to the mouth while the area remains anaesthetised.
\end{itemize}

\section*{Prognosis and Outcomes}
The prognosis for the vast majority of patients receiving local anaesthetics is excellent, with the numbness resolving spontaneously and completely within a few hours, depending on the agent and dose used. Lignocaine typically wears off within 1 to 2 hours, whereas long-acting agents like bupivacaine can last up to 8 hours. When systemic toxicity occurs, the prompt administration of lipid emulsion therapy has dramatically improved outcomes, with most patients recovering fully without long-term sequelae. Permanent nerve injury is extremely rare, with an incidence of less than 1 in 10,000 cases.

\section*{Prevention and Self-Care}
Precautionary steps and post-procedure care are key to preventing injuries and adverse outcomes associated with local anaesthesia.

\begin{itemize}
    \item \textbf{Medical disclosure:} Inform the healthcare provider of all medical conditions (especially liver, kidney, or heart disease), current medications, pregnancy status, and any past reactions to anaesthetics.
    \item \textbf{Protect the anaesthetised area:} Avoid touching, scratching, or rubbing the numb area to prevent accidental skin damage.
    \item \textbf{Dietary precautions:} Do not consume hot food or drinks, or chew solid food, until oral sensation has fully returned to prevent burning or biting the mouth.
    \item \textbf{Activity restriction:} Avoid strenuous exercise or placing weight on a limb that is numb or weak from a regional nerve block until full motor control returns.
    \item \textbf{Wound protection:} Keep dressing sites clean and dry, and follow specific post-procedure wound care instructions provided by the clinician.
\end{itemize}

\section*{Sources and Further Reading}
The following organisations provide reputable and comprehensive information on local anaesthetics, regional anaesthesia, and patient safety.

\begin{itemize}
    \item Australian and New Zealand College of Anaesthetists. \textit{Patient information on anaesthesia and pain management.} https://www.anzca.edu.au/
    \item Healthdirect Australia. \textit{Local anaesthetics: uses, risks, and recovery.} https://www.healthdirect.gov.au/local-anaesthetics
    \item NHS. \textit{Overview of local anaesthesia and minor procedures.} https://www.nhs.uk/
    \item Mayo Clinic. \textit{Local and regional anaesthesia patient guides.} https://www.mayoclinic.org/
    \item National Institute for Health and Care Excellence. \textit{Guidelines on local and regional anaesthetic safety.} https://www.nice.org.uk/
\end{itemize}
